\section{Literature Review}\label{sec:literature}

The theory this thesis builds on splits into a few strands that have largely
developed in parallel, and it is worth seeing where each of them stands before
combining them. The first strand is the classical correspondence between regular
expressions and finite automata. The standard constructions that realise an
expression as a machine are Thompson's translation into a nondeterministic automaton
with epsilon transitions~\cite{Thompson1968}, the position automaton of
Glushkov~\cite{Glushkov1962}, and the state-graph method of McNaughton and
Yamada~\cite{McNaughtonYamada1960}. These remain the textbook routes from syntax to
machine, and they treat the automaton as an object to be built alongside the
expression rather than read off from it.

A second strand reasons about regular expressions algebraically. Conway's account of
regular algebra and finite machines~\cite{C71} organises the equational theory,
while Salomaa gave complete axiom systems for the algebra of regular
events~\cite{Salomaa} and Kozen later established completeness for Kleene algebra
and the algebra of regular events~\cite{Kozen}. This line is concerned with when two
expressions denote the same language and with axiomatising that equivalence, which
is a different question from how to match efficiently, but it is closely related,
since deciding equivalence is one of the things a good normal form for expressions
should help with.

The derivative approach sits between these two and is the one this thesis takes as
its starting point. Brzozowski introduced derivatives of regular
expressions~\cite{B62} and showed that differentiating an expression with respect to
a symbol, together with a test for whether the empty word is accepted, is enough to
both decide membership and construct a deterministic automaton, provided expressions
are identified up to a few simplifications so that only finitely many derivatives
arise. Owens, Reppy and Turon re-examined the construction from a modern
implementation perspective~\cite{OwensReppyTuron} and showed that it remains both
elegant and practical, particularly because it extends smoothly to intersection and
complement. Derivatives are attractive precisely because the automaton is latent in
the expression rather than assembled separately, which is also what makes the method
a good fit for mechanised reasoning.

That fit has been borne out by work on formalising regular languages in proof
assistants. The development this thesis extends descends from the Coquand--Siles
decision procedure for regular-expression equivalence in type
theory~\cite{CoquandRegexEquality} and its accompanying Rocq
formalisation of Brzozowski derivatives~\cite{RocqRegexpBrzozowski}, carried out in
the Rocq proof assistant~\cite{Rocq}. Mechanising this material forces every
simplification and every finiteness argument to be made explicit, which is exactly
where paper proofs tend to be informal, and it is the foundation on which the
similarity layer of the present work is added.

The fuzzy side of the picture has its own substantial literature, though it has
mostly grown from a different root. The dominant tradition equips automata and
expressions with membership values rather than crisp acceptance, so that a word is
recognised to a degree. Li and Pedrycz studied fuzzy finite automata and fuzzy
regular expressions with membership values in lattice-ordered
monoids~\cite{LP06}, Stamenkovi\'c and \'Ciri\'c gave constructions of fuzzy automata
from fuzzy regular expressions~\cite{SC12}, and Garhwal, Jiwari and Tomasiello
surveyed fuzzy regular languages and fuzzy automata more broadly~\cite{GJT17}. A
related quantitative and coalgebraic perspective appears in the quantitative Kleene
coalgebras of Silva, Bonchi, Bonsangue and Rutten~\cite{SilvaBonchiBonsangueRutten2011}
and, more recently, in the extension of the quantitative pattern-matching paradigm by
Alves, Kesner and Ramos~\cite{AlvesKesnerRamos2025}. The technical backdrop for all
of this is the theory of triangular norms and fuzzy logic, of which Metcalfe's notes
give a compact account~\cite{met05}.

What is distinctive about the route taken here is the source of the fuzziness. Rather
than attaching weights to transitions or to acceptance, this thesis keeps acceptance
crisp and instead loosens the matching of symbols, so that a symbol matches any other
symbol that is similar to it within a chosen threshold. This is the same move that
Sessa made in approximate logic programming, where unification is replaced by
matching up to a similarity relation~\cite{Sessa2002}, transposed from resolution to
derivatives. The payoff is that, under the minimum T-norm, the fuzzy theory reduces
to the crisp one over a quotient alphabet, which keeps the derivative machinery and
its formalisation intact instead of requiring a separate weighted theory. The
present work therefore connects the derivative-based and formally verified treatment
of regular expressions with the similarity-based style of approximate matching, and
the membership-value tradition cited above is the natural point of comparison for
what such a similarity-based account does and does not capture.
